\documentclass[12pt,a4paper]{article}

% ===================== PACKAGES =====================
\usepackage[a4paper,margin=1in]{geometry}
\usepackage{amsmath, amssymb}
\usepackage{setspace}
\usepackage{enumitem}
\usepackage{titlesec}
\usepackage{fancyhdr}
\usepackage{hyperref}

% ===================== FORMATTING =====================
\setstretch{1.2}
\setlist[itemize]{noitemsep}
\setlist[enumerate]{noitemsep}

\pagestyle{fancy}
\fancyhf{}
\rhead{CSE400}
\lhead{Lecture 4 Scribe}
\cfoot{\thepage}

\titleformat{\section}{\large\bfseries}{\thesection.}{0.5em}{}
\titleformat{\subsection}{\normalsize\bfseries}{\thesubsection}{0.5em}{}

% ===================== DOCUMENT =====================
\begin{document}

\begin{center}
    {\Large \textbf{CSE400 -- Fundamentals of Probability in Computing}}\\[4pt]
    {\large Lecture 4: Joint Probability and Conditional Probability}\\[6pt]
    \textit{Prepared for Exam Revision}
\end{center}

\vspace{1em}

% =====================================================
\section{Introduction to Probability Theory}

This lecture establishes the foundational framework required to define joint probability and conditional probability. The development proceeds through precise definitions, probability axioms, and their consequences.

% =====================================================
\section{Experiments, Outcomes, and Events}

\subsection{Experiment}

\textbf{Definition:}  
An \textit{experiment} $(E)$ is a procedure that we perform which produces some result.

\textbf{Example:}  
Tossing a coin five times, denoted as $E_5$.

\textbf{Logical Role:}  
An experiment defines the probabilistic setting. Probabilities are assigned only after the experiment is specified.

% -----------------------------------------------------
\subsection{Outcome}

\textbf{Definition:}  
An \textit{outcome} $(\xi)$ is a possible result of an experiment.

\textbf{Example:}
\[
\xi_1 = \text{HHTHT}
\]

\textbf{Logical Role:}  
Each outcome represents one complete realization of the experiment.

% -----------------------------------------------------
\subsection{Event}

\textbf{Definition:}  
An \textit{event} is a certain set of outcomes of an experiment.

\textbf{Example:}  
Let $C$ be the event in experiment $E_5$ defined as:
\[
C = \{\text{all outcomes consisting of an even number of heads}\}
\]

\textbf{Reasoning:}  
An event is a subset of the sample space and groups outcomes based on a condition of interest.

% =====================================================
\section{Sample Space}

\subsection{Definition}

The \textit{sample space} $(S)$ is the collection of all possible distinct outcomes of an experiment.

\textbf{Properties:}
\begin{itemize}
    \item Mutually Exclusive: Only one outcome can occur at a time.
    \item Collectively Exhaustive: No possible outcome lies outside the sample space.
\end{itemize}

% -----------------------------------------------------
\subsection{Types of Sample Spaces}

The sample space may be:
\begin{itemize}
    \item Discrete
    \item Countably Infinite
    \item Continuous
\end{itemize}

\textbf{Examples listed:}
\begin{itemize}
    \item Flipping a fair coin once
    \item Rolling a cubical die
    \item Rolling two dice
    \item Flipping a coin until tails occurs
    \item Random number generator in $[0,1)$
\end{itemize}

The structure of the sample space determines how probability axioms are applied.

% =====================================================
\section{Probability and Axioms}

\subsection{Definition of Probability}

Probability is:
\begin{itemize}
    \item A measure of the likelihood of various events, or
    \item A function assigning a numerical value to an event.
\end{itemize}

% -----------------------------------------------------
\subsection{Probability Axioms}

The axioms are taken to be self-evident.

For any event $A$:

\begin{enumerate}
    \item \textbf{Boundedness}
    \[
    0 \leq \Pr(A) \leq 1
    \]

    \item \textbf{Certain Event}
    \[
    \Pr(S) = 1
    \]

    \item \textbf{Additivity for Mutually Exclusive Events}
    
    If $A \cap B = \varnothing$, then
    \[
    \Pr(A \cup B) = \Pr(A) + \Pr(B)
    \]

    \item \textbf{Countable Additivity}
    \[
    \Pr\left(\bigcup_{i=1}^{\infty} A_i\right)
    =
    \sum_{i=1}^{\infty} \Pr(A_i)
    \]
\end{enumerate}

These axioms define the mathematical behavior of probability.

% =====================================================
\section{Corollary}

\subsection{Corollary 2.1}

For a finite number $M$ of mutually exclusive events:
\[
\Pr\left(\bigcup_{i=1}^{M} A_i\right)
=
\sum_{i=1}^{M} \Pr(A_i)
\]

Axiom 3 reduces to this corollary when the sample space is finite.

% =====================================================
\section{Propositions}

\subsection{Proposition 2.1 (Complement Rule)}

\[
\Pr(A^c) = 1 - \Pr(A)
\]

\textbf{Reasoning:}
\begin{enumerate}
    \item $A$ and $A^c$ are mutually exclusive.
    \item $A \cup A^c = S$.
    \item Since $\Pr(S) = 1$, subtracting $\Pr(A)$ gives $\Pr(A^c)$.
\end{enumerate}

% -----------------------------------------------------
\subsection{Proposition 2.2 (Monotonicity)}

If $A \subset B$, then
\[
\Pr(A) \leq \Pr(B)
\]

Since $B$ contains all outcomes of $A$, its probability cannot be smaller.

% =====================================================
\section{Joint Probability}

\subsection{Definition}

The joint probability of events $A$ and $B$ is:
\[
\Pr(A \cap B)
\]

It represents the probability that both $A$ and $B$ occur simultaneously.

\subsection{Examples Mentioned}

\begin{itemize}
    \item Card Deck Example
    \item Costume Party Example
\end{itemize}

In each case:
\begin{enumerate}
    \item Define individual events.
    \item Identify their intersection.
    \item Assign probability to the intersection.
\end{enumerate}

% =====================================================
\section{Conditional Probability}

\subsection{Definition}

For $\Pr(B) > 0$:
\[
\Pr(A \mid B) = \frac{\Pr(A \cap B)}{\Pr(B)}
\]

\subsection{Step-by-Step Interpretation}

\begin{enumerate}
    \item Occurrence of $B$ restricts the sample space to $B$.
    \item Favorable outcomes are in $A \cap B$.
    \item Division by $\Pr(B)$ normalizes probability in the reduced space.
\end{enumerate}

\subsection{Examples Mentioned}

\begin{itemize}
    \item Cards without replacement
    \item Game of Poker
    \item The Missing Key
\end{itemize}

Each example:
\begin{enumerate}
    \item Identifies the conditioning event.
    \item Restricts the sample space.
    \item Applies the conditional probability formula.
\end{enumerate}

% =====================================================
\section{Key Exam Points}

\begin{itemize}
    \item Clearly distinguish experiment, outcome, event, and sample space.
    \item State axioms before using them.
    \item Joint probability corresponds to intersection.
    \item Conditional probability is normalized joint probability.
    \item Ensure $\Pr(B) > 0$ before computing $\Pr(A \mid B)$.
\end{itemize}

\end{document}
